% TODO
% Response to referee 2, comment 3
Plan:
\begin{enumerate}
\item All of our theory indeed suggests that interplane coupling plays little role here.
So, you could say that the van der Waals-nature of the material studied here does not matter for the Higgs-mode observation.
\item However, the point is that van der Waals materials are both technologically quite interesting for feasibility reasons, and they are just a new family of materials from which new phenomena are frequently being discovered.
Therefore, it feels relevant to mention that the discovery occurred in a van der Waals material, even though it is not relevant to the mechanism of the phenomenon we discovered.
\item We can also potentially cook up some reason van der Waals is important?
For example, the 1-dimensional chain motif is certainly important here.
Does the same motif occur in non-van der Waals materials?
\textcolor{red}{Ajesh said he would think about it}
\item For the question about ferroelectric sliding modes, Ajesh said he is thinking about it!
\end{enumerate}

In our work, the physical mechanism for ferroelectricity is magneto-electric coupling to a spin-spiral magnetic order, and as such, the van der Waals (vdW) nature of the material does not play a direct role in this regard.
As the referee points out, in contrast, there are other mechanisms for how ferroelectricity could arise in vdW materials.
In sliding ferroelectrics, such as parallel-stacked bilayers of hexagonal boron nitride or rhombohedral transition-metal dichalcogenides, the monolayers themselves lack inversion symmetry.
An inter-layer charge transfer can this occur resulting in an out-of-plane electric polarization. 
What are the collective modes of such systems? 
As layers slide relative to each other the polarization gets modified, and the associated sliding modes are referred to as ferroelectric sliding modes. 

This mechanism is absent in \ce{CuBr2}. Each \ce{CuBr2} layer is itself
centrosymmetric, with an inversion center at every Cu site. Any two layers related by a translation, including an arbitrary relative sliding, therefore remain related by inversion through the midpoint between them, so interlayer sliding cannot induce a polarization. 
Both the mechanism of ferroelectricity and the nature of the collective modes in \ce{CuBr2} are therefore distinct from those of sliding ferroelectrics. This is supported by our data: both modes soften as the temperature approaches the magnetic transition at $T_N$ and vanish above it, as expected for collective modes of the magnetic order rather than for lattice vibrations.
